% Do not change document class, margins, fonts, etc.
\documentclass[a4paper,oneside,bibliography=totoc]{scrartcl}

% some useful packages (add more as needed)
\usepackage[utf8]{inputenc}
\usepackage{graphicx}
\usepackage{latexsym}
\usepackage{amsmath}
\usepackage{amssymb}
\usepackage{tabularx}
\usepackage{booktabs}
\usepackage{algorithm} % you can modify the algorithm style to your liking
\usepackage{algorithmic}
\usepackage{csquotes}
\renewcommand{\algorithmiccomment}[1]{\hfill\textit{// #1}}
\usepackage[usenames,dvipsnames]{xcolor}
\usepackage[colorlinks,citecolor=Green]{hyperref} % you may change/remove the colors
\usepackage{lipsum} % you do not need this
\usepackage{environ}
\usepackage{numprint}
\usepackage{tikz}
\usetikzlibrary{positioning,shapes.gates.logic.US,shapes.gates.logic.IEC,automata,arrows, backgrounds}
\usepackage{natbib}
\bibliographystyle{chicagoa}
\setcitestyle{authoryear,round,semicolon,aysep={},yysep={,}} \let\cite\citep

% example environments and commands (add more as needed)
\newtheorem{definition}{Definition} \newtheorem{proposition}{Proposition}

\newcommand\todo[1]{\textcolor{red}{\framebox{\footnotesize TODO}\expandafter\ifx\expandafter\relax\detokenize{#1}\relax\else\ #1\fi}}


% DO NOT MODIFY
\newlength\oldparskip
\NewEnviron{answer}[1][8cm]{%
  \sbox{0}{\parbox[t][#1][t]{0.99\linewidth}{%
      \setlength{\oldparskip}{\parskip}%
      \setlength{\parskip}{.5em}%
      \BODY%
      ~%
      \setlength{\parskip}{\oldparskip}%
    }}%

  \begin{tikzpicture}
    % inner sep restores the \fbox padding
    \node[inner sep=\fboxsep, outer sep=0pt] (A) {\usebox{0}};
    % border thickness like \fboxrule
    \draw[line width=\fboxrule] (A.south west) rectangle (A.north east);
  \end{tikzpicture}%
  \vskip\baselineskip
}


\begin{document}

\title{DL26, Assignment 2}
\author{\todo{Student 1 name (login, matriculation number)}\\
	\todo{Student 2 name (login, matriculation number)}}
\date{\today}

\maketitle

%%%%%%%%%%%%%%%%%%%%%%%%%%%%%%%%%%%%%%%%%%%%%%%%%%%%%%%%%%%%%%%%%%%%%%%%%%%%%%%%
% Task 2. %%%%%%%%%%%%%%%%%%%%%%%%%%%%%%%%%%%%%%%%%%%%%%%%%%%%%%%%%%%%%%%%%%%%%%
%%%%%%%%%%%%%%%%%%%%%%%%%%%%%%%%%%%%%%%%%%%%%%%%%%%%%%%%%%%%%%%%%%%%%%%%%%%%%%%%

\section*{Task 2a)}
\begin{answer}[8cm]
	% YOUR ANSWER HERE
\end{answer}

\section*{Task 2c)}
\begin{answer}[16cm]
	% YOUR ANSWER HERE
\end{answer}

\section*{Task 2d): i), ii), and iii)}
\begin{answer}[12cm]
	% YOUR ANSWER HERE
\end{answer}

\section*{Task 2e)}
\begin{answer}[7cm]
	% YOUR ANSWER HERE
\end{answer}

\section*{Task 2f)}
\begin{answer}[4cm]
	% YOUR ANSWER HERE
\end{answer}

\section*{Task 2h)}
\begin{answer}[4cm]
	% YOUR ANSWER HERE
\end{answer}

\section*{Task 2i)}
\begin{answer}[6cm]
	% YOUR ANSWER HERE
\end{answer}

%%%%%%%%%%%%%%%%%%%%%%%%%%%%%%%%%%%%%%%%%%%%%%%%%%%%%%%%%%%%%%%%%%%%%%%%%%%%%%%%
% Task 3. %%%%%%%%%%%%%%%%%%%%%%%%%%%%%%%%%%%%%%%%%%%%%%%%%%%%%%%%%%%%%%%%%%%%%%
%%%%%%%%%%%%%%%%%%%%%%%%%%%%%%%%%%%%%%%%%%%%%%%%%%%%%%%%%%%%%%%%%%%%%%%%%%%%%%%%

\section*{Task 3c)}
\begin{answer}[4cm]
	% YOUR ANSWER HERE
\end{answer}

\section*{Task 3d)}
\begin{answer}[19cm]
	% YOUR ANSWER HERE
\end{answer}

\section*{Task 3e)}
\begin{answer}[19cm]
	% YOUR ANSWER HERE
\end{answer}

\section*{Task 3f)}
\begin{answer}[19cm]
	% YOUR ANSWER HERE
\end{answer}

%%%%%%%%%%%%%%%%%%%%%%%%%%%%%%%%%%%%%%%%%%%%%%%%%%%%%%%%%%%%%%%%%%%%%%%%%%%%%%%%
% Task 4. %%%%%%%%%%%%%%%%%%%%%%%%%%%%%%%%%%%%%%%%%%%%%%%%%%%%%%%%%%%%%%%%%%%%%%
%%%%%%%%%%%%%%%%%%%%%%%%%%%%%%%%%%%%%%%%%%%%%%%%%%%%%%%%%%%%%%%%%%%%%%%%%%%%%%%%

\section*{Task 4)}
\begin{answer}[19cm]
	% YOUR ANSWER HERE
\end{answer}

%%%%%%%%%%%%%%%%%%%%%%%%%%%%%%%%%%%%%%%%%%%%%%%%%%%%%%%%%%%%%%%%%%%%%%%%%%%%%%%%
% AI Declaration. %%%%%%%%%%%%%%%%%%%%%%%%%%%%%%%%%%%%%%%%%%%%%%%%%%%%%%%%%%%%%%
%%%%%%%%%%%%%%%%%%%%%%%%%%%%%%%%%%%%%%%%%%%%%%%%%%%%%%%%%%%%%%%%%%%%%%%%%%%%%%%%

\clearpage
\section*{Ehrenwörtliche Erklärung}

Ich versichere, dass ich die beiliegende Bachelor-, Master-, Seminar-, oder
Projektarbeit ohne Hilfe Dritter und ohne Benutzung anderer als der angegebenen
Quellen und in der untenstehenden Tabelle angegebenen Hilfsmittel angefertigt
und die den benutzten Quellen wörtlich oder inhaltlich entnommenen Stellen als
solche kenntlich gemacht habe. Diese Arbeit hat in gleicher oder ähnlicher Form
noch keiner Prüfungsbehörde vorgelegen. Ich bin mir bewusst, dass eine falsche
Erklärung rechtliche Folgen haben wird.

% Declare below which AI tools you used in the process of writing your work,
% including text, image, code, and data generation.
%
% The AI declaration must include at least one an entry for every graded subtask
% (i.e., where points are specified). The entry should state whether or not you
% used AI to solve parts of the task. Using AI for things such as language
% improvements or API search does not count as "Solved with AI", but things such
% as code gen, writing drafts, or solving math tasks does.
%
% We may grade tasks that you (partially or fully) solved with AI more
% stringently. So make sure that whenever you use AI, the results are correct.
% Our recommendation is to not use AI to solve the tasks, as this undermines the
% learning goals of this course and, ultimately, your own education.
\begin{center}
  \textbf{Declaration of Used AI Tools} \\[.3em]
  \begin{tabularx}{\textwidth}{lclXc}
    \toprule
    Task & Solved with AI? & Tool & Purpose & Useful? \\
    \midrule
    1b & Yes & ChatGPT & Generate draft solution & yes \\
    1c \\
    1d \\
    2a \\
    2b \\
    2c \\
    2d \\
    3a \\
    3b \\
    \bottomrule
  \end{tabularx}
\end{center}

\vspace{2cm}
\noindent Unterschrift\\
\noindent Mannheim, den XX.XX.2025 \hfill


\end{document}
